\begin{theorem}[Characters of the norm quotient]
\label{mod:ramification-normcharacters}
Let $K/F$ be a finite extension of nonarchimedean local fields and write
\[
 Q_{K/F}=F^\times/N_{K/F}(K^\times)
 =F^\times/\operatorname{range}(N_{K/F}).
\]
The quotient is always oriented with $F^\times$ as numerator.  Lean defines
$\widehat Q_{K/F}$ to be the continuous homomorphisms
$Q_{K/F}\to\mathbf C^\times$, and equips the project's type
\[
 S(K/F)=\{\mu:F^\times\to\mathbf C^\times:
             \mu\circ N_{K/F}=1\}
\]
with its pointwise commutative-group structure.  Quotient lifting and
pullback are inverse multiplicative equivalences
\[
 \widehat Q_{K/F}\simeq S(K/F).
\]
Thus a character is trivial on the norm subgroup if and only if it factors
through $Q_{K/F}$; the construction uses the quotient topology and makes
representative independence explicit.  When the norm range is open, every
algebraic character of $Q_{K/F}$ is continuous, and complex finite-group
duality gives
\[
 \#S(K/F)=\#Q_{K/F}.
\]

Now assume that $K/F$ is prime cyclic, let
$\ell=[K:F]$, and first consider the ramified case.  Supply a lower break
$t$, the equality $f(K/F)=1$, an integral uniformizer $\pi_K$, and the
monogenicity equality required by
module~\ref{mod:ramification-normfiltration}.  Put
\[
 C_t=\operatorname{gr}^tU_F/
   \operatorname{range}(\overline N_t).
\]
Lean constructs the surjection $U_F^t\to C_t$ and proves that its kernel is
the honest critical norm image.  The global intersection and join formulas
from the norm-filtration theorem then give
\[
 \#Q_{K/F}
 =\#\bigl(U_F^t/(N(K^\times)\cap U_F^t)\bigr)
 =\#C_t=\ell.
\]
In particular, this cardinality is derived from the proved norm-filtration
API and does not use local class field theory.  The norm quotient and the
complete character group are finite, and
\[
 \#Q_{K/F}=\#S(K/F)=\ell.
\]
The trivial character is included among these $\ell$ characters.  The
subtype of nontrivial characters has cardinality $\ell-1$.

The exact conductor statements are as follows.
\begin{itemize}
\item The trivial norm character has multiplicative conductor $0$.
\item If $K/F$ is unramified, every norm character, including every
  nontrivial one, has conductor $0$.
\item In the ramified case, every norm character is trivial on
  $U_F^{t+1}$.  A norm character trivial already on $U_F^t$ is globally
  trivial, by the global join formula.  Hence every nontrivial norm
  character has exact conductor $t+1$.
\item Consequently the tame case $t=0$ gives exact conductor $1$, while
  the wild case $t>0$ gives exact conductor $t+1\ge2$.
\end{itemize}
All these are proofs of \texttt{IsMultiplicativeConductor}; no conductor is
assumed as a field of a norm-character structure.

In the unramified case Lean reuses the explicit normalized-order quotient
from module~\ref{mod:ramification-unramifiedcompatibility}.  It re-exports
\[
 \#Q_{K/F}=\#S(K/F)=[K:F]
\]
and the corresponding finiteness results.  For an unramified prime cyclic
extension it constructs both an arbitrary cyclic enumeration
\[
 \operatorname{Multiplicative}(\mathbf Z/\ell\mathbf Z)
 \simeq S(K/F)
\]
and the enumeration aligned with a chosen uniformizer $\pi_F$,
\[
 S(K/F)\simeq\mu_\ell(\mathbf C),\qquad
 \mu\longmapsto\mu(\pi_F).
\]
The latter is an equivalence, so every $\ell$-th root occurs uniquely as a
uniformizer value.

Combining the two branches, Lean exports the global non-instance theorem
\texttt{primeCyclicNormCharacter\_finite}.  It classifies a prime-cyclic
extension according to whether $e(K/F)=1$.  The unramified branch uses the
preceding unramified finiteness theorem.  In the ramified branch, the green
preparation package from
module~\ref{mod:ramification-primecyclicpreparation} supplies the canonical
lower break, residue-degree equality, integral uniformizer, and monogenicity
proof required by the ramified finiteness theorem.  Consequently
\[
 \operatorname{Finite}(S(K/F))
\]
holds for every prime-cyclic extension, still including the trivial norm
character.  Downstream files install this result locally with
\texttt{letI := primeCyclicNormCharacter\_finite ...}.

For both ramified and unramified prime-degree cases, Lean exports explicit
finite sets of all norm characters and all cyclic indices, cyclic
equivalences, the fact that index zero is the trivial character, an
equivalence between nonzero indices and nontrivial characters, and
full-product reindexing formulas.  These products retain the zero/trivial
term.  Multiplication by any fixed norm character is also exported as a
permutation, together with invariance of a full product under that
reindexing.

Finally, let $\chi$ be a quasi-character carrying an exact conductor proof,
and suppose a depth $T$ is supplied at which every norm character is
trivial.  For every $\mu\in S(K/F)$, Lean constructs the actual least
conductor of $\mu\chi$ by minimizing the unit layers on which it is trivial.
It then chooses
\[
 \mu_0\in S(K/F)
\]
whose twist has least conductor in the orbit.  The chosen twist is packaged
with its exact conductor proof, and Lean proves both
\[
 m(\mu_0\chi)\le m(\nu\chi)
 \qquad(\nu\in S(K/F))
\]
and the orbit-stable formulation: whenever $m$ is an exact conductor of
$\nu(\mu_0\chi)$,
\[
 m(\mu_0\chi)\le m.
\]
No uniqueness of the minimizing character is asserted.

\LeanFile{LanglandsFirstMainLemma/Ramification/NormCharacters.lean}
\lean{LanglandsFirstMainLemma.NormQuotient}
\lean{LanglandsFirstMainLemma.NormQuotientCharacter}
\lean{LanglandsFirstMainLemma.normQuotientCharacterEquivNormCharacter}
\lean{LanglandsFirstMainLemma.normQuotientCharacterEquivMonoidHom}
\lean{LanglandsFirstMainLemma.normQuotientCharacter_card}
\lean{LanglandsFirstMainLemma.normCharacter_card_of_isOpen}
\lean{LanglandsFirstMainLemma.CriticalNormCokernel}
\lean{LanglandsFirstMainLemma.criticalNormCokernelProjection}
\lean{LanglandsFirstMainLemma.criticalNormCokernelProjection_ker}
\lean{LanglandsFirstMainLemma.criticalNormTargetQuotientEquivCokernel}
\lean{LanglandsFirstMainLemma.normQuotient_card}
\lean{LanglandsFirstMainLemma.normQuotient_finite}
\lean{LanglandsFirstMainLemma.ramifiedNormCharacter_card}
\lean{LanglandsFirstMainLemma.ramifiedNormCharacter_finite}
\lean{LanglandsFirstMainLemma.unramifiedNormQuotient_card}
\lean{LanglandsFirstMainLemma.unramifiedNormCharacter_card}
\lean{LanglandsFirstMainLemma.unramifiedNormCharacter_finite}
\lean{LanglandsFirstMainLemma.primeCyclicNormCharacter_finite}
\lean{LanglandsFirstMainLemma.trivialNormCharacter_conductor}
\lean{LanglandsFirstMainLemma.ramifiedNormCharacter_conductor}
\lean{LanglandsFirstMainLemma.tameNormCharacter_conductor}
\lean{LanglandsFirstMainLemma.wildNormCharacter_conductor}
\lean{LanglandsFirstMainLemma.unramifiedNormCharacter_conductor}
\lean{LanglandsFirstMainLemma.normCharacterFintype}
\lean{LanglandsFirstMainLemma.normCharacterFinset}
\lean{LanglandsFirstMainLemma.normCharacter_prod_mul}
\lean{LanglandsFirstMainLemma.nontrivialNormCharacterIndexEquiv}
\lean{LanglandsFirstMainLemma.ramifiedNormCharacterZModEquiv}
\lean{LanglandsFirstMainLemma.ramifiedNontrivialNormCharacter_card}
\lean{LanglandsFirstMainLemma.ramifiedNormCharacterProduct_eq_zmodProduct}
\lean{LanglandsFirstMainLemma.unramifiedNormCharacterZModEquiv}
\lean{LanglandsFirstMainLemma.unramifiedNormCharacterEvaluationEquiv}
\lean{LanglandsFirstMainLemma.existsUnique_normCharacter_uniformizer_value}
\lean{LanglandsFirstMainLemma.unramifiedNontrivialNormCharacter_card}
\lean{LanglandsFirstMainLemma.unramifiedNormCharacterProduct_eq_zmodProduct}
\lean{LanglandsFirstMainLemma.normCharacterTwistData}
\lean{LanglandsFirstMainLemma.minimalConductorTwist}
\lean{LanglandsFirstMainLemma.minimalConductorTwistData}
\lean{LanglandsFirstMainLemma.minimalConductorTwist_isConductor}
\lean{LanglandsFirstMainLemma.minimalConductorTwist_minimal}
\lean{LanglandsFirstMainLemma.minimalConductorTwist_actualMinimal}
\leanok
\SourceText{Corollary~\ref{cor:norm-characters} and
Lemma~\ref{lem:norm-character-conversion}.}
\uses{mod:ramification-normfiltration,mod:delta-firstmainstatement,
mod:ramification-unramifiedcompatibility,
mod:ramification-primecyclicpreparation}
\FormalizationContract
The quotient/type conversion proved here is the equivalence between
characters of $F^\times/N(K^\times)$ and the project's
\texttt{NormCharacter}.  The manuscript's
Lemma~\ref{lem:norm-character-conversion} is a later stationary
norm--trace phase identity; this node does not claim that separate identity.

The ramified cardinality proof uses only the critical cokernel cardinality,
critical successor image, and global intersection/join fields of
module~\ref{mod:ramification-normfiltration}.  It assumes no local
class-field-theory identification.  Positive ramified conductor statements
explicitly assume that the character is nontrivial; the trivial character
is retained in every finite type, enumeration, and full product.  The
unramified and ramified conductor-zero boundaries are kept separate.
\end{theorem}
